\section{El formato GWB: la base de trabajo}
\label{sec:A02GWB}


El cartucho de tipo GWB (Game \& Watch Base) es el cartucho de PARTIDA de la autoría: contiene sobre todo la imagen de partida con los objetos dibujados, todavía sin
segmentos identificados. Se \emph{entrega como una CARPETA} (no un ZIP) con estos ficheros en su interior:

\begin{itemize}
	\item Un \texttt{.json} con los datos del cartucho GWB (cualquier nombre; si hay
	varios, la prioridad del apartado anterior).
	
	\item \texttt{imagenes.img\_f}: OPCIONAL. La imagen de fondo sin los segmentos. Hay 
	cartuchos sin arte de fondo. Si \texttt{img\_f} está declarada, el 	fichero referenciado
	DEBE estar presente en la carpeta; sin fondo, el simulador pinta la pantalla en blanco.
	
	\item \texttt{imagenes.img\_s}: OBLIGATORIO. La imagen con todos los segmentos encendidos pero sin 
	el fondo, OBLIGATORIA siempre (de ella se extrae la geometría): el fichero DEBE estar presente en la carpeta. Si existe \texttt{img\_f}
	debe medir exactamente lo mismo; si no es así hay que dar un aviso por consola NO bloqueante.
	
	\item \texttt{imagenes.img\_m}: OPCIONAL. La portada de la máquina, que se usará en la
	pantalla de selección de juegos como referencia gráfica del cartucho a cargar. Si se declara, el fichero referenciado DEBE estar presente.
\end{itemize}

Los campos del JSON asociado a un GWB:

\begin{table}[H]
	\centering
	\begin{tabular}{llp{10cm}}
		\toprule
		Campo & Tipo & Contenido \\
		\midrule
		\texttt{tipo} & cadena & \texttt{"GWB"} \\
		\texttt{version} & número & Versión del formato (actualmente 1) \\
		\texttt{nombre} & cadena & Nombre del cartucho \\
		\texttt{estilo} & cadena & Estilo de la máquina (debe coincidir con el setup) \\
		\texttt{imagenes} & objeto & Tres campos \texttt{\{img\_m, img\_f, img\_s\}} con los  nombres de los ficheros de imagen \\
		\texttt{resolucion} & objeto & OPCIONAL: \texttt{\{ancho, alto\}} en píxeles de la pantalla LCD. Si este dato no se indica usará la definida en el GWS cargado. \\
		\texttt{botones} & array & OPCIONAL: Los nombres que damos a los botones. \texttt{[\{numero, id\}]} por cada botón \\
		\texttt{memoria} & objeto & OPCIONAL: \texttt{\{RAM, FLASH\}} en KB, para el porcentaje de ocupaci\u00f3n al generar el GWC (conceptos comunes y \ref{sec:F5I20}) \\
		\bottomrule
	\end{tabular}
	\caption{Campos del json del GWB}
	\label{tab:A02GWB}
\end{table}

Un GWB NO lleva sprites, segmentos ni autoría: al cargarlo, el editor
descompone \texttt{img\_s} automáticamente (extracción de componentes
negras) y el diseñador trabaja a partir de ahí. 

El campo \texttt{numero} de un botón es el número físico del hardware (los
\texttt{id} son etiquetas legibles como \texttt{IZQ} o \texttt{Start}).

\begin{lstlisting}[
	language=json, 
	frame=single, 
	breaklines=true, 
	basicstyle=\ttfamily\small,
	caption={Ejemplo de json para un cartucho GWB},
	label={lst:A02GWB}
	]
	{
		"tipo": "GWB",
		"version": 1,
		"nombre": "Parachute",
		"estilo": "2",
		"imagenes": {
			"img_m": "parachute_m.png",
			"img_f": "parachute_f.png",
			"img_s": "parachute_s.png"
		},
		"botones": [
		{"numero": 0, "id": "IZQ"},
		{"numero": 4, "id": "DER"},
		{"numero": "A", "id": "Start"}
		]
	}
\end{lstlisting}
